\documentclass[a4paper]{article}
\usepackage{parskip}
\usepackage{lipsum}
% \usepackage{newtxtext}

\def\nterm {Fall}
\def\nyear {2026}
\def\ncourse {Mathematical Studies Analysis}

\makeatletter
\ifx \nauthor\undefined
  \def\nauthor{Runqiu Ye}
\else
\fi

\author{Notes taken by \nauthor\vspace{5pt}\\
  Carnegie Mellon University}
\date{\nterm\ \nyear}

\usepackage{alltt}
\usepackage{amsfonts}
\usepackage{amsmath}
\usepackage{amssymb}
\usepackage{amsthm}
\usepackage{booktabs}
\usepackage{caption}
\usepackage{enumitem}
\usepackage{fancyhdr}
\usepackage{graphicx}
\usepackage{mathdots}
\usepackage{mathtools}
\usepackage{microtype}
\usepackage{multirow}
\usepackage{pdflscape}
\usepackage{pgfplots}
\usepackage{siunitx}
\usepackage{slashed}
\usepackage{tabularx}
\usepackage{tikz}
\usepackage{tkz-euclide}
\usepackage[normalem]{ulem}
\usepackage[all]{xy}
\usepackage{imakeidx}
\usepackage[includehead,includefoot,heightrounded,left=1.25in,right=1.25in]{geometry}
\usepackage{mathrsfs}
\usepackage{stmaryrd}
\usepackage{bbm}

\makeindex[intoc, title=Index]
\indexsetup{othercode={\lhead{\emph{Index}}}}

\ifx \nextra \undefined
  \usepackage[pdftex,
    hidelinks,
    pdfauthor={Dexter Chua},
    pdfsubject={\ncourse},
    pdftitle={\ncourse},
    pdfkeywords={Cambridge Mathematics Maths Math \nterm\ \nyear\ \ncourse}]{hyperref}
  \title{\ncourse}
\else
  \usepackage[pdftex,
    hidelinks,
    pdfauthor={Dexter Chua},
    pdfsubject={Cambridge Maths Notes: \ncourse\ (\nextra)},
    pdftitle={\ncourse\ (\nextra)},
    pdfkeywords={Cambridge Mathematics Maths Math \nterm\ \nyear\ \ncourse\ \nextra}]{hyperref}

  \title{\ncourse \\ {\Large \nextra}}
  \renewcommand\printindex{}
\fi

\pgfplotsset{compat=1.12}

\pagestyle{fancyplain}
\ifx \ncoursehead \undefined
  \def\ncoursehead{\ncourse}
\fi

\lhead{\emph{\nouppercase{\leftmark}}}
\ifx \nextra \undefined
  \rhead{
    \ifnum\thepage=1
    \else
      \ncoursehead
    \fi}
\else
  \rhead{
    \ifnum\thepage=1
    \else
      \ncoursehead \ (\nextra)
    \fi}
\fi
\usetikzlibrary{arrows.meta}
\usetikzlibrary{decorations.markings}
\usetikzlibrary{decorations.pathmorphing}
\usetikzlibrary{positioning}
\usetikzlibrary{fadings}
\usetikzlibrary{intersections}
\usetikzlibrary{cd}

\newcommand*{\Cdot}{{\raisebox{-0.25ex}{\scalebox{1.5}{$\cdot$}}}}
% \newcommand {\pd}[2][ ]{
%   \ifx #1 { }
%     \frac{\partial}{\partial #2}
%   \else
%     \frac{\partial^{#1}}{\partial #2^{#1}}
%   \fi
% }
\newcommand{\pd}[2]{\frac{\partial #1}{\partial #2}}
\ifx \nhtml \undefined
\else
  \renewcommand\printindex{}
  \DisableLigatures[f]{family = *}
  \let\Contentsline\contentsline
  \renewcommand\contentsline[3]{\Contentsline{#1}{#2}{}}
  \renewcommand{\@dotsep}{10000}
  \newlength\currentparindent
  \setlength\currentparindent\parindent

  \newcommand\@minipagerestore{\setlength{\parindent}{\currentparindent}}
  \usepackage[active,tightpage,pdftex]{preview}
  \renewcommand{\PreviewBorder}{0.1cm}

  \newenvironment{stretchpage}%
  {\begin{preview}\begin{minipage}{\hsize}}%
  {\end{minipage}\end{preview}}
  \AtBeginDocument{\begin{stretchpage}}
  \AtEndDocument{\end{stretchpage}}

  \newcommand{\@@newpage}{\end{stretchpage}\begin{stretchpage}}

  \let\@real@section\section
  \renewcommand{\section}{\@@newpage\@real@section}
  \let\@real@subsection\subsection
  \renewcommand{\subsection}{\@ifstar{\@real@subsection*}{\@@newpage\@real@subsection}}
\fi
\ifx \ntrim \undefined
\else
  \usepackage{geometry}
  \geometry{
    papersize={379pt, 699pt},
    textwidth=345pt,
    textheight=596pt,
    left=17pt,
    top=54pt,
    right=17pt
  }
\fi

\ifx \nisofficial \undefined
  \let\@real@maketitle\maketitle
  \renewcommand{\maketitle}{\@real@maketitle}
\else
\fi

% Theorems
\theoremstyle{definition}
\newtheorem*{aim}{Aim}
\newtheorem*{axiom}{Axiom}
\newtheorem*{claim}{Claim}
\newtheorem*{cor}{Corollary}
\newtheorem*{conjecture}{Conjecture}
\newtheorem*{defi}{Definition}
\newtheorem*{eg}{Example}
\newtheorem*{ex}{Exercise}
\newtheorem*{sol}{Solution}
\newtheorem*{fact}{Fact}
\newtheorem*{law}{Law}
\newtheorem*{lemma}{Lemma}
\newtheorem*{notation}{Notation}
\newtheorem*{prop}{Proposition}
\newtheorem*{question}{Question}
\newtheorem*{rrule}{Rule}
\newtheorem*{thm}{Theorem}
\newtheorem*{assumption}{Assumption}

\newtheorem*{remark}{Remark}
\newtheorem*{warning}{Warning}
\newtheorem*{exercise}{Exercise}

\newtheorem{nthm}{Theorem}[section]
\newtheorem{nlemma}[nthm]{Lemma}
\newtheorem{nprop}[nthm]{Proposition}
\newtheorem{ncor}[nthm]{Corollary}

\renewcommand{\labelitemi}{--}
\renewcommand{\labelitemii}{$\circ$}

\let\stdsection\section
\renewcommand\section{\newpage\stdsection}

% Strike through
\def\st{\bgroup \ULdepth=-.55ex \ULset}

%%%%%%%%%%%%%%%%%%%%%%%%%
%%%%% Maths Symbols %%%%%
%%%%%%%%%%%%%%%%%%%%%%%%%
\newcommand{\zseqinfl}[1]{\{ #1 \}_{l=0}^\infty}
\newcommand{\zseqinfj}[1]{\{ #1 \}_{j=0}^\infty}
\newcommand{\zseqinfi}[1]{\{ #1 \}_{i=0}^\infty}
\newcommand{\zseqinfn}[1]{\{ #1 \}_{n=0}^\infty}
\newcommand{\zseqinfk}[1]{\{ #1 \}_{k=0}^\infty}
\newcommand{\zseqinfm}[1]{\{ #1 \}_{m=0}^\infty}

\newcommand{\seqinfl}[1]{\{ #1 \}_{l=1}^\infty}
\newcommand{\cupinfj}{\bigcup_{j=1}^\infty}
\newcommand{\capinfj}{\bigcap_{j=1}^\infty}
\newcommand{\suminfj}{\sum_{j=1}^\infty}
\newcommand{\seqinfj}[1]{\{ #1 \}_{j=1}^\infty}
\newcommand{\cupinfi}{\bigcup_{i=1}^\infty}
\newcommand{\capinfi}{\bigcap_{i=1}^\infty}
\newcommand{\suminfi}{\sum_{i=1}^\infty}
\newcommand{\seqinfi}[1]{\{ #1 \}_{i=1}^\infty}
\newcommand{\cupinfn}{\bigcup_{n=1}^\infty}
\newcommand{\capinfn}{\bigcap_{n=1}^\infty}
\newcommand{\suminfn}{\sum_{n=1}^\infty}
\newcommand{\seqinfn}[1]{\{ #1 \}_{n=1}^\infty}
\newcommand{\cupinfk}{\bigcup_{k=1}^\infty}
\newcommand{\sqcupinfk}{\bigsqcup_{k=1}^\infty}
\newcommand{\capinfk}{\bigcap_{k=1}^\infty}
\newcommand{\suminfk}{\sum_{k=1}^\infty}
\newcommand{\seqinfk}[1]{\{ #1 \}_{k=1}^\infty}
\newcommand{\cupinfm}{\bigcup_{m=1}^\infty}
\newcommand{\capinfm}{\bigcap_{m=1}^\infty}
\newcommand{\suminfm}{\sum_{m=1}^\infty}
\newcommand{\seqinfm}[1]{\{ #1 \}_{m=1}^\infty}
\newcommand{\limsupn}{\limsup_{n \to \infty}}
\newcommand{\liminfn}{\liminf_{n \to \infty}}

\newcommand{\cupj}{\bigcup_{j=1}}
\newcommand{\capj}{\bigcap_{j=1}}
\newcommand{\sumj}{\sum_{j=1}}
\newcommand{\seqj}[1]{\{ #1 \}_{j=1}}
\newcommand{\cupi}{\bigcup_{i=1}}
\newcommand{\capi}{\bigcap_{i=1}}
\newcommand{\sumi}{\sum_{i=1}}
\newcommand{\seqi}[1]{\{ #1 \}_{i=1}}
\newcommand{\cupn}{\bigcup_{n=1}}
\newcommand{\capn}{\bigcap_{n=1}}
\newcommand{\sumn}{\sum_{n=1}}
\newcommand{\seqn}[1]{\{ #1 \}_{n=1}}
\newcommand{\cupk}{\bigcup_{k=1}}
\newcommand{\capk}{\bigcap_{k=1}}
\newcommand{\sumk}{\sum_{k=1}}
\newcommand{\seqk}[1]{\{ #1 \}_{k=1}}
\newcommand{\cupm}{\bigcup_{m=1}}
\newcommand{\capm}{\bigcap_{m=1}}
\newcommand{\summ}{\sum_{m=1}}
\newcommand{\seqm}[1]{\{ #1 \}_{m=1}}

% mathfrak 
\newcommand{\afk}{\mathfrak{A}}
\newcommand{\bfk}{\mathfrak{B}}
\newcommand{\hfk}{\mathfrak{H}}

% mathcal 
\newcommand{\acal}{\mathcal{A}}
\newcommand{\bcal}{\mathcal{B}}
\newcommand{\ccal}{\mathcal{C}}
\newcommand{\ecal}{\mathcal{E}}
\newcommand{\fcal}{\mathcal{F}}
\newcommand{\gcal}{\mathcal{G}}
\newcommand{\hcal}{\mathcal{H}}
\newcommand{\ncal}{\mathcal{N}}
\newcommand{\ocal}{\mathcal{O}}
\newcommand{\rcal}{\mathcal{R}}
\newcommand{\tcal}{\mathcal{T}}
\newcommand{\ucal}{\mathcal{U}}
\def\L{\mathcal{L}}
\def\H{\mathcal{H}}
\renewcommand{\P}{\mathcal{P}}

% Analysis
\newcommand{\om}{\mu^*}
\newcommand{\fs}{F_\sigma}
\newcommand{\gd}{G_\delta}
\newcommand{\mf}{\mathfrak{M}}
\newcommand{\nf}{\mathfrak{N}}
\newcommand{\Sfin}{S_{\text{fin}}}

\newcommand{\loc}{\mathrm{loc}}
\DeclareMathOperator{\leb}{Leb}
\DeclareMathOperator{\Rec}{Rec}
\DeclareMathOperator{\Cube}{Cube}
\DeclareMathOperator{\RCube}{RCube}
\newcommand{\RCuber}{\RCube_r}
\DeclareMathOperator{\Reco}{Rec^{\circ}}
\DeclareMathOperator{\Cubeo}{Cube^{\circ}}
\DeclareMathOperator{\RCubeo}{RCube^{\circ}}
\newcommand{\RCubeor}{\RCube^{\circ}_r}

\DeclareMathOperator{\Var}{Var}

% Matrix groups
\newcommand{\GL}{\mathrm{GL}}
\newcommand{\Or}{\mathrm{O}}
\newcommand{\PGL}{\mathrm{PGL}}
\newcommand{\PSL}{\mathrm{PSL}}
\newcommand{\PSO}{\mathrm{PSO}}
\newcommand{\PSU}{\mathrm{PSU}}
\newcommand{\SL}{\mathrm{SL}}
\newcommand{\SO}{\mathrm{SO}}
\newcommand{\Spin}{\mathrm{Spin}}
\newcommand{\Sp}{\mathrm{Sp}}
\newcommand{\SU}{\mathrm{SU}}
\newcommand{\U}{\mathrm{U}}
\newcommand{\Mat}{\mathrm{Mat}}

% Matrix algebras
\newcommand{\gl}{\mathfrak{gl}}
\newcommand{\ort}{\mathfrak{o}}
\newcommand{\so}{\mathfrak{so}}
\newcommand{\su}{\mathfrak{su}}
\newcommand{\uu}{\mathfrak{u}}
\renewcommand{\sl}{\mathfrak{sl}}

% Special sets
\newcommand{\C}{\mathbb{C}}
\newcommand{\CP}{\mathbb{CP}}
\newcommand{\GG}{\mathbb{G}}
\newcommand{\N}{\mathbb{N}}
\newcommand{\Q}{\mathbb{Q}}
\newcommand{\R}{\mathbb{R}}
\newcommand{\RP}{\mathbb{RP}}
\newcommand{\T}{\mathbb{T}}
\newcommand{\Z}{\mathbb{Z}}
% \renewcommand{\H}{\mathbb{H}}

% Brackets
\renewcommand{\bar}[1]{\overline{#1}}
\renewcommand{\tilde}[1]{\widetilde{#1}}
\renewcommand{\hat}[1]{\widehat{#1}}
\newcommand{\floor}[1]{\left\lfloor #1 \right\rfloor}
\newcommand{\ceil}[1]{\left\lceil #1 \right\rceil}
\newcommand{\abs}[1]{\left\lvert #1\right\rvert}
\newcommand{\bket}[1]{\left\lvert #1\right\rangle}
\newcommand{\brak}[1]{\left\langle #1 \right\rvert}
\newcommand{\braket}[1]{\left\langle #1 \right\rangle}
\newcommand{\bra}{\left\langle}
\newcommand{\ket}{\right\rangle}
\newcommand{\norm}[1]{\left\lVert #1\right\rVert}
\newcommand{\normalorder}[1]{\mathop{:}\nolimits\!#1\!\mathop{:}\nolimits}
\newcommand{\tv}[1]{|#1|}
\newcommand{\bbracket}[1]{\left\llbracket #1 \right\rrbracket}
\renewcommand{\vec}[1]{\boldsymbol{\mathbf{#1}}}
\DeclareMathOperator{\curl}{curl}
\DeclareMathOperator{\diverge}{div}
\DeclareMathOperator{\dist}{dist}

% not-math
\newcommand{\bolds}[1]{{\bfseries #1}}
\newcommand{\cat}[1]{\mathsf{#1}}
\newcommand{\ph}{\,\cdot\,}
\newcommand{\term}[1]{\emph{#1}\index{#1}}
\newcommand{\phantomeq}{\hphantom{{}={}}}

% Probability
\DeclareMathOperator{\ind}{\mathbf{1}}
\DeclareMathOperator{\Bernoulli}{Bernoulli}
\DeclareMathOperator{\betaD}{beta}
\DeclareMathOperator{\bias}{bias}
\DeclareMathOperator{\binomial}{binomial}
\DeclareMathOperator{\corr}{corr}
\DeclareMathOperator{\cov}{cov}
\DeclareMathOperator{\gammaD}{gamma}
\DeclareMathOperator{\mse}{mse}
\DeclareMathOperator{\multinomial}{multinomial}
\DeclareMathOperator{\Poisson}{Poisson}
\DeclareMathOperator{\Exp}{Exp}
\DeclareMathOperator{\unif}{unif}
\DeclareMathOperator{\var}{var}
\newcommand{\E}{\mathbb{E}}
\renewcommand{\Pr}{\mathbb{P}}

% Algebra
\DeclareMathOperator{\adj}{adj}
\DeclareMathOperator{\Ann}{Ann}
\DeclareMathOperator{\Aut}{Aut}
\DeclareMathOperator{\Char}{char}
\DeclareMathOperator{\disc}{disc}
\DeclareMathOperator{\dom}{dom}
\DeclareMathOperator{\fix}{fix}
\DeclareMathOperator{\Hom}{Hom}
\DeclareMathOperator{\id}{id}
\DeclareMathOperator{\image}{image}
\DeclareMathOperator{\im}{im}
\DeclareMathOperator{\tr}{tr}
\DeclareMathOperator{\Tr}{Tr}
\newcommand{\Bilin}{\mathrm{Bilin}}
\newcommand{\Frob}{\mathrm{Frob}}

% Others
\newcommand\ad{\mathrm{ad}}
\newcommand\Art{\mathrm{Art}}
\newcommand{\B}{\mathcal{B}}
\newcommand{\cU}{\mathcal{U}}
\newcommand{\Der}{\mathrm{Der}}
\newcommand{\D}{\mathrm{D}}
\newcommand{\dR}{\mathrm{dR}}
\newcommand{\exterior}{\mathchoice{{\textstyle\bigwedge}}{{\bigwedge}}{{\textstyle\wedge}}{{\scriptstyle\wedge}}}
\newcommand{\F}{\mathbb{F}}
\newcommand{\G}{\mathcal{G}}
\newcommand{\Gr}{\mathrm{Gr}}
\newcommand{\haut}{\mathrm{ht}}
\newcommand{\Hol}{\mathrm{Hol}}
\newcommand{\hol}{\mathfrak{hol}}
\newcommand{\Id}{\mathrm{Id}}
\newcommand{\lie}[1]{\mathfrak{#1}}
\newcommand{\op}{\mathrm{op}}
\newcommand{\Oc}{\mathcal{O}}
\newcommand{\pr}{\mathrm{pr}}
\newcommand{\Ps}{\mathcal{P}}
\newcommand{\pt}{\mathrm{pt}}
\newcommand{\qeq}{\mathrel{``{=}"}}
\newcommand{\Rs}{\mathcal{R}}
\newcommand{\Vect}{\mathrm{Vect}}
\newcommand{\wsto}{\stackrel{\mathrm{w}^*}{\to}}
\newcommand{\wt}{\mathrm{wt}}
\newcommand{\wto}{\stackrel{\mathrm{w}}{\to}}
% \renewcommand{\d}{\; \mathrm{d}}
\renewcommand{\d}{\, d}
\newcommand{\dd}{\partial}
\newcommand{\nsubset}{\not\subset}
\renewcommand{\epsilon}{\varepsilon}
\renewcommand{\phi}{\varphi}

\let\Im\relax
\let\Re\relax

\DeclareMathOperator{\RHS}{RHS}
\DeclareMathOperator{\LHS}{LHS}
\DeclareMathOperator{\step}{Step}
\DeclareMathOperator{\reg}{Reg}
\DeclareMathOperator{\ran}{range}
\DeclareMathOperator{\area}{area}
\DeclareMathOperator{\card}{card}
\DeclareMathOperator{\ccl}{ccl}
\DeclareMathOperator{\ch}{ch}
\DeclareMathOperator{\cl}{cl}
\DeclareMathOperator{\cls}{\overline{\mathrm{span}}}
\DeclareMathOperator{\coker}{coker}
\DeclareMathOperator{\conv}{conv}
\DeclareMathOperator{\cosec}{cosec}
\DeclareMathOperator{\cosech}{cosech}
\DeclareMathOperator{\covol}{covol}
\DeclareMathOperator{\diag}{diag}
\DeclareMathOperator{\diam}{diam}
\DeclareMathOperator{\Diff}{Diff}
\DeclareMathOperator{\End}{End}
\DeclareMathOperator{\energy}{energy}
\DeclareMathOperator{\erfc}{erfc}
\DeclareMathOperator{\erf}{erf}
\DeclareMathOperator*{\esssup}{ess\,sup}
\DeclareMathOperator{\ev}{ev}
\DeclareMathOperator{\Ext}{Ext}
\DeclareMathOperator{\fst}{fst}
\DeclareMathOperator{\Fit}{Fit}
\DeclareMathOperator{\Frac}{Frac}
\DeclareMathOperator{\Gal}{Gal}
\DeclareMathOperator{\gr}{gr}
\DeclareMathOperator{\hcf}{hcf}
\DeclareMathOperator{\Im}{Im}
\DeclareMathOperator{\Ind}{Ind}
\DeclareMathOperator{\Int}{Int}
\DeclareMathOperator{\Isom}{Isom}
\DeclareMathOperator{\lcm}{lcm}
\DeclareMathOperator{\length}{length}
\DeclareMathOperator{\Lie}{Lie}
\DeclareMathOperator{\like}{like}
\DeclareMathOperator{\Lk}{Lk}
\DeclareMathOperator{\Maps}{Maps}
\DeclareMathOperator{\orb}{orb}
\DeclareMathOperator{\ord}{ord}
\DeclareMathOperator{\otp}{otp}
\DeclareMathOperator{\poly}{poly}
\DeclareMathOperator{\rank}{rank}
\DeclareMathOperator{\rel}{rel}
\DeclareMathOperator{\Rad}{Rad}
\DeclareMathOperator{\Re}{Re}
\DeclareMathOperator*{\res}{res}
\DeclareMathOperator{\Res}{Res}
\DeclareMathOperator{\Ric}{Ric}
\DeclareMathOperator{\rk}{rk}
\DeclareMathOperator{\Rees}{Rees}
\DeclareMathOperator{\Root}{Root}
\DeclareMathOperator{\sech}{sech}
\DeclareMathOperator{\sgn}{sgn}
\DeclareMathOperator{\snd}{snd}
\DeclareMathOperator{\Spec}{Spec}
\DeclareMathOperator{\spn}{span}
\DeclareMathOperator{\stab}{stab}
\DeclareMathOperator{\St}{St}
\DeclareMathOperator{\supp}{supp}
\DeclareMathOperator{\spt}{spt}
\DeclareMathOperator{\Syl}{Syl}
\DeclareMathOperator{\Sym}{Sym}
\DeclareMathOperator{\vol}{vol}

\pgfarrowsdeclarecombine{twolatex'}{twolatex'}{latex'}{latex'}{latex'}{latex'}
\tikzset{->/.style = {decoration={markings,
          mark=at position 1 with {\arrow[scale=2]{latex'}}},
      postaction={decorate}}}
\tikzset{<-/.style = {decoration={markings,
          mark=at position 0 with {\arrowreversed[scale=2]{latex'}}},
      postaction={decorate}}}
\tikzset{<->/.style = {decoration={markings,
          mark=at position 0 with {\arrowreversed[scale=2]{latex'}},
          mark=at position 1 with {\arrow[scale=2]{latex'}}},
      postaction={decorate}}}
\tikzset{->-/.style = {decoration={markings,
          mark=at position #1 with {\arrow[scale=2]{latex'}}},
      postaction={decorate}}}
\tikzset{-<-/.style = {decoration={markings,
          mark=at position #1 with {\arrowreversed[scale=2]{latex'}}},
      postaction={decorate}}}
\tikzset{->>/.style = {decoration={markings,
          mark=at position 1 with {\arrow[scale=2]{latex'}}},
      postaction={decorate}}}
\tikzset{<<-/.style = {decoration={markings,
          mark=at position 0 with {\arrowreversed[scale=2]{twolatex'}}},
      postaction={decorate}}}
\tikzset{<<->>/.style = {decoration={markings,
          mark=at position 0 with {\arrowreversed[scale=2]{twolatex'}},
          mark=at position 1 with {\arrow[scale=2]{twolatex'}}},
      postaction={decorate}}}
\tikzset{->>-/.style = {decoration={markings,
          mark=at position #1 with {\arrow[scale=2]{twolatex'}}},
      postaction={decorate}}}
\tikzset{-<<-/.style = {decoration={markings,
          mark=at position #1 with {\arrowreversed[scale=2]{twolatex'}}},
      postaction={decorate}}}

\tikzset{circ/.style = {fill, circle, inner sep = 0, minimum size = 3}}
\tikzset{scirc/.style = {fill, circle, inner sep = 0, minimum size = 1.5}}
\tikzset{mstate/.style={circle, draw, blue, text=black, minimum width=0.7cm}}

\tikzset{eqpic/.style={baseline={([yshift=-.5ex]current bounding box.center)}}}
\tikzset{commutative diagrams/.cd,cdmap/.style={/tikz/column 1/.append style={anchor=base east},/tikz/column 2/.append style={anchor=base west},row sep=tiny}}

\definecolor{mblue}{rgb}{0.2, 0.3, 0.8}
\definecolor{morange}{rgb}{1, 0.5, 0}
\definecolor{mgreen}{rgb}{0.1, 0.4, 0.2}
\definecolor{mred}{rgb}{0.5, 0, 0}

\def\drawcirculararc(#1,#2)(#3,#4)(#5,#6){%
  \pgfmathsetmacro\cA{(#1*#1+#2*#2-#3*#3-#4*#4)/2}%
  \pgfmathsetmacro\cB{(#1*#1+#2*#2-#5*#5-#6*#6)/2}%
  \pgfmathsetmacro\cy{(\cB*(#1-#3)-\cA*(#1-#5))/%
    ((#2-#6)*(#1-#3)-(#2-#4)*(#1-#5))}%
  \pgfmathsetmacro\cx{(\cA-\cy*(#2-#4))/(#1-#3)}%
  \pgfmathsetmacro\cr{sqrt((#1-\cx)*(#1-\cx)+(#2-\cy)*(#2-\cy))}%
  \pgfmathsetmacro\cA{atan2(#2-\cy,#1-\cx)}%
  \pgfmathsetmacro\cB{atan2(#6-\cy,#5-\cx)}%
  \pgfmathparse{\cB<\cA}%
  \ifnum\pgfmathresult=1
    \pgfmathsetmacro\cB{\cB+360}%
  \fi
  \draw (#1,#2) arc (\cA:\cB:\cr);%
}
\newcommand\getCoord[3]{\newdimen{#1}\newdimen{#2}\pgfextractx{#1}{\pgfpointanchor{#3}{center}}\pgfextracty{#2}{\pgfpointanchor{#3}{center}}}

\newcommand\qedshift{\vspace{-17pt}}
\newcommand\fakeqed{\pushQED{\qed}\qedhere}

\def\Xint#1{\mathchoice
  {\XXint\displaystyle\textstyle{#1}}%
  {\XXint\textstyle\scriptstyle{#1}}%
  {\XXint\scriptstyle\scriptscriptstyle{#1}}%
  {\XXint\scriptscriptstyle\scriptscriptstyle{#1}}%
  \!\int}
\def\XXint#1#2#3{{\setbox0=\hbox{$#1{#2#3}{\int}$}
    \vcenter{\hbox{$#2#3$}}\kern-.5\wd0}}
\def\ddashint{\Xint=}
\def\dashint{\Xint-}

\newcommand\separator{{\centering\rule{2cm}{0.2pt}\vspace{2pt}\par}}

\newenvironment{own}{\color{gray!70!black}}{}

\newcommand\makecenter[1]{\raisebox{-0.5\height}{#1}}

\mathchardef\mdash="2D

\newenvironment{significant}{\begin{center}\begin{minipage}{0.9\textwidth}\centering\em}{\end{minipage}\end{center}}
\DeclareRobustCommand{\rvdots}{%
  \vbox{
    \baselineskip4\p@\lineskiplimit\z@
    \kern-\p@
    \hbox{.}\hbox{.}\hbox{.}
  }}
\DeclareRobustCommand\tph[3]{{\texorpdfstring{#1}{#2}}}
\makeatother


\author{Notes taken by \nauthor\vspace{5pt}\\ 
Lectures by Giovanni Leoni\vspace{5pt}\\
Carnegie Mellon University}

\newcommand{\TODO}{\textcolor{red}{\textbf{*** TO-DO ***}}}

\begin{document}
\maketitle

\tableofcontents

\section{The fundamental theorem of calculus}

\subsection{Differentiation of monotone functions}

\subsubsection*{Continuity}

\begin{defi}
  Let $X$ be a set and $g : X \to [0, \infty)$. Define infinite sum
  \[
    \sum_{x \in X} g(x) \defeq \left\{ \sum_{x \in Y} g(x) :
    Y \subset X \text{ finite} \right\}.
  \]
\end{defi}

\begin{thm}
  Let $X$ be a set and $g : X \to [0, \infty)$. If $\sum_{x \in X} g(x) <
    \infty$, then $\{ x \in X : g(x) > 0\}$ is countable.
\end{thm}

\begin{proof}
  Let $M = \sum_{x \in X} g(x)$ and $X_k = \{ x \in X : g(x) > k^{-1} \}$
  for $k \in \N$. Let $Y \subset X_k$ finite, then
  \[
    M \geq \sum_{x \in Y} g(x) > k^{-1} \abs{Y}.
  \]
  Hence $\abs{Y} \leq k M$ and thus $\abs{X_k} \leq k M$. It follows that
  $\{ x \in X : g(x) > 0 \} = \cupinfk X_k$ is countable.
\end{proof}

\begin{thm}
  Let $I \subset \R$ be an interval and $f : I \to \R$ an increasing
  function. Then, $f$ is continuous at all but countably many $x \in I$.
\end{thm}

\begin{proof}
  First assume $I = [a, b]$. Then for $x \in (a, b)$, there exists
  $\lim_{y \to x^-} f(y) = f(x-)$ and $\lim_{y \to x^+} f(y) = f(x+)$.
  Since the function is increasing, $f(x+) \geq f(x-)$. Now we can define
  $g (x) = f(x+) - f(x-)$. It is clear that $\sum_{x \in [a, b]} g(x) \leq
    f(b) - f(a)$. By the previous theorem, $f(x+) = f(x-)$ for all but
  countably many $x \in [a, b]$.

  The interval $I \subset \R$ can be written as a countable union of
  finite closed intervals. The proof is then complete.
\end{proof}

\subsubsection*{Differentiation}

\begin{thm}[Lebesgue's differentiation theorem]
  Let $I \subset \R$ be an interval and let $f : I \to \R$ be a monotone
  function. Then $f$ is differentiable a.e.\ on $I$.
\end{thm}

To show this, we need a few lemmas.

\begin{lemma}
  Let $I \subset \R$ be an interval, let $f : I \to \R$, and let
  \[
    E = \left\{ x \in I^\circ : \text{ there exist $f'_+ (x)$ and
  $f'_- (x)$ and $f'_+(x) \neq f'_-(x)$} \right\}.
  \]
  Then $E$ is countable.
\end{lemma}

\begin{proof}
  First observe that for any function $g: I \to \R$, the set of its strict
  local minimum points in $I^\circ$ is countable. Set
  \[
    E_- = \left\{ x \in E : f'_-(x) < f'_+ (x) \right\}.
  \]
  For every $r \in \Q$, define
  \[
    E_r = \left\{ x \in E_- : f'_-(x) < r < f'_+ (x) \right\},
  \]
  and let $g_r(t) = f(t) - rt$ for $t \in I$. By density of rationals,
  $E_- = \bigcup_{r \in \Q} E_r$. If $x \in E_r$, then
  \[
    (g_r)'_- (x) < 0 < (g_r)'_+ (x).
  \]
  It follows that $g_r(y) > g_r(x)$ for every $y \neq x$ sufficiently
  close to $x$, i.e.\ $x$ is a strict local minimum point of $g_r$.
  Therefore $E_r$ is countable and thus $E_-$ is countable. Applying the
  same argument to $-f$ completes the proof for the lemma.
\end{proof}

\begin{lemma}
  Let $E \subset \R$ be bounded, and $\fcal$ be a family of open intervals
  such that each $x \in E$ is the left endpoint of an interval $(x, x +
    h_x) \in \fcal$. Then there for every $\epsilon > 0$, there exist
  disjoint intervals $I_1, \dots, I_n \in \fcal$ such that
  \[
    \L_o^1 \left( E \cap \cupk^n I_k \right)  \geq \L_o^1 (E) - \epsilon.
  \]
\end{lemma}

\begin{proof}
  If $E = \emptyset$, take the empty collection. If $E \neq \emptyset$ and
  $\L_o^1 (E) \leq \epsilon$, any one interval in $\fcal$ suffices. Thus
  assume that $L_o^1 (E) > \epsilon$. For each $x \in E$, choose an
  interval $(x, x + h_x) \in \fcal$ and define
  \[
    E_n = \left\{ x \in E : h_x > n^{-1} \right\}.
  \]
  Note that $E_n \subset E_{n + 1}$ and $E = \cupinfn E_n$. By coninuity
  from below of Lebesgue outer measure (exercise), we have $\L_o^1 (E_n)
  \to \L_o^1 (E)$. Hence we can find $m$ large enough such that $\L_o^1
  (E_m) > \L_o^1 (E) - \frac{\epsilon}{2}$. For each $k \in \N$, let
  $\eta_k = \epsilon 2^{- k - 1}$, then $\suminfk \eta_k =
  \frac{\epsilon}{2}$.

  Let $a_1 = \inf E_m$. By definition of infimum we can find $x_1 \in E_m$
  with $a_1 \leq x_1 < a_1 + \eta_1$, and correspondingly $I_1 = (x_1 +
  h_1) \in \fcal$ with $h_1 > m^{-1}$. Suppose $I_1, \dots, I_k$ has been
  chosen. If $\left\{ x \in E_m : x \geq x_k + h_k \right\} = \emptyset$,
  we stop. Otherwise let $a_{k + 1} = \inf \left\{ x \in E_m : x \geq x_k
  + h_k \right\}$. Then there is $x_{k + 1} \in E_m$ such that $a_{k + 1}
  \leq x_{k + 1} \leq a_{k + 1} + \eta_{k + 1}$ and correspondingly $I_{k
  + 1} = (x_{k + 1}, x_{k + 1} + h_{k + 1}) \in \fcal$ with $h_{k + 1} >
  m^{-1}$. Note that
  \[
    x_{k + 1} \geq x_k + h_k > x_k + m^{-1} ,
  \]
  and $E_m \subset E$ is bounded, the construction must stop after
  finitely many steps. Moreover, $I_1, \dots, I_n$ are disjoint by
  construction.

  Now let
  \[
    S = \cupk^n I_n, \quad T = \cupk^n (x_k - \eta_k, x_k].
  \]
  The points of $E_m$ skipped at the $k$-th step belong to $[a_k, x_k]
  \subset (x_k - \eta_k, x_k]$. Thus $E_m \subset S \cup T$, then
  \begin{align*}
    \L_o^1 (E) - \frac{\epsilon}{2}
     & < \L_o^1 (E_m) \leq \L_o^1 (E_m \cap S) + \L_o^1 (E_m \cap T) \\
     & \leq \L_o^1 (E_m \cap S) + \sumk^n \eta_k
    \leq \L_o^1 (E_m \cap S) + \frac{\epsilon}{2}.
  \end{align*}
  This concludes the proof.
\end{proof}

\begin{lemma}
  Let $E \subset \R$ be bounded, $\fcal$ be family of open intervals such
  that for all $x \in E$ and $\delta > 0$, there exists $(x, x + h_x) \in \fcal$
  with $0 \leq h_x < \delta$. Then for all $\epsilon > 0$, there exists
  disjoint intervals $I_1, \dots, I_m \in \fcal$ such that
  \[
    \L_o^1 \left( E \cap \cupi^m I_i \right) \geq \L_o^1(E) - \epsilon,
    \quad
    \sumi^m \L^1 (I_i) \leq \L_o^1 (E) + \epsilon.
  \]
\end{lemma}

\begin{proof}
  By the definiton of Lebesgue outer measure, find an bounded open set $U
    \supset E$ such that $\L_o^1 (U) \leq \L_o^1 (E) + \epsilon$. Let
  $\fcal'$ be the subfamily of intervals $(x, x + h_x) \in \fcal$
  contained in $U$. Note that for each $x \in E \subset U$, there must
  exists at least one such interval, since $U$ contains a ball centered at
  $x$ and there are intervals of arbitrarily small length.

  Apply the previous lemma to the family $\fcal'$ to find disjoint
  intervals $I_1, \dots, I_n \subset \fcal'$ such that
  \[
    \L_o^1 \left( E \cap \cupk^n I_k \right)  \geq \L_o^1 (E) - \epsilon.
  \]
  Since the intervals are disjoint and contined in $U$, it follows that
  \[
    \sumk^n \L^1 (I_k)  \leq \L_o^1 (U) \leq \L_o^1 (E) + \epsilon.
  \]
  The proof is then complete.
\end{proof}

\begin{defi}
  For a function $f : E \to \R$, define
  \[
    \liminf_{x \to x_0} f(x) = \sup_{r > 0} \inf_{0 < d(x, x_0) < r} f(x),
    \quad
    \limsup_{x \to x_0} f(x) = \inf_{r > 0} \sup_{0 < d(x, x_0) < r} f(x),
  \]
\end{defi}

\begin{ex}
  $\liminf_{x \to x_0} f(x) = l \in \R$ iff for every
  $\epsilon > 0$:
  \begin{enumerate}
    \item there exists $\delta > 0$ such that $l - \epsilon < f(x)$ for
          all $x \in E$ with $0 < d(x, x_0) < \delta$.
    \item For every $\delta > 0$, there exists $x \in E$ with $0 < d(x,
          x_0) < \delta$ such that $f(x) < l + \epsilon$.
  \end{enumerate}

  Also, $\liminf_{x \to x_0} f(x) = -\infty$ iff for every $M > 0$ and
  $\delta > 0$, there exists $x \in E$ with $0 < d(x, x_0) < \delta$ such
  that $f(x) < - M$.
\end{ex}

Now we are ready to show Lebesgue's differentiation theorem. Recall the
statement of the theorem below:

\begin{thm}[Lebesgue's differentiation theorem]
  Let $I \subset \R$ be an interval and let $f : I \to \R$ be a monotone
  function. Then $f$ is differentiable a.e.\ on $I$.
\end{thm}

\begin{proof}
  Define \textbf{Dini's derivatives} as follows:
  \begin{align*}
    \bar{D}_+ f(x_0)       & = \limsup_{x \to x_0^+} \frac{f(x) - f(x_0)}{x - x_0}, &
    \underline{D}_+ f(x_0) & = \liminf_{x \to x_0^+} \frac{f(x) - f(x_0)}{x - x_0},   \\
    \bar{D}_- f(x_0)       & = \limsup_{x \to x_0^-} \frac{f(x) - f(x_0)}{x - x_0}, &
    \underline{D}_- f(x_0) & = \liminf_{x \to x_0^-} \frac{f(x) - f(x_0)}{x - x_0}.   \\
  \end{align*}

  WLOG assume $f$ is increasing, and assume $I = [a, b]$. Let
  \begin{align*}
    E & = \left\{ x \in I^\circ : \underline{D}_+ f(x) < \bar{D}_+ f(x) \right\}   \\
      & = \bigcup_{\substack{r, s \in \Q \\ 0 < r < s}}
    \left\{ x \in I^\circ : \underline{D}_+ f(x) < r < s < \bar{D}_+ f(x) \right\} \\
      & = \bigcup_{\substack{r, s \in \Q \\ 0 < r < s}} E_{r, s}
  \end{align*}
  It suffices to show that $\L^1 (E_{r, s}) = 0$. Let $\epsilon > 0$ be given.

  Let $x \in E_{r, s}$. We have
  \[
    \liminf_{h \to 0^+} \frac{f(x + h) - f(x)}{h} < r.
  \]
  Then for each $\delta > 0$, there exists $0 < h_x < \delta$ such that
  $h_x^{-1} (f(x + h_x) - f(x)) < r$. Let $\fcal$ be the collection of
  $(x, x + h_x)$, where $h_x$ is selected as aforementioned while $\delta$
  ranges over all positive value, and $x$ ranges over $E_{r, s}$. It is
  then clear that $\fcal$ satisfies the condition for the previous lemma.

  Apply previous lemma, there exists $I_1, \dots, I_m \in \fcal$ disjoint
  such that
  \begin{align*}
     & \L_o^1 \left( E_{r, s} \cap \cupi^m I_i \right) \geq \L_o^1 (E_{r, s}) - \epsilon, \tag{1} \\
     & \sumi^m \L^1 (I_i) \leq \L_o^1 (E_{r,s}) + \epsilon. \tag{2}
  \end{align*}
  Write $I_i = (x_i, x_i + h_i)$, then
  \begin{align*}
    \sumi^n f(x_i + h_i) - f(x_i)
    < r \sumi^n h_i
    \leq r \L_o^1 (E_{r,s}) + r \epsilon. \tag{3}
  \end{align*}

  Now let $V = \cupi^m (x_i, x_i + h_i)$, which is open. If $x \in E_{r,
  s} \cap V$, we have
  \[
    s < \limsup_{t \to 0} \frac{f(x + t) - f(x)}{t}.
  \]
  Then for every $\delta > 0$, there exists $f_x$ such that $0 < t_x <
  \delta$ and $s < t_x^{-1} (f(x + t_x) - f(x))$. Define $\gcal$ to be the
  collection of $(x, x + t_x)$ that is contained in $V$, where $t_x$ is
  selected as aforementioned while $\delta > 0$ ranges over all positive
  value, and $x$ ranges over $E_{r, s} \cap V$.

  The previous lemma again gives $J_1, \dots, J_l \in \gcal$ disjoint such
  that
  \begin{align*}
    \L_o^1 \left( E_{r, s} \cap V \cap \cupk^l J_k \right)
    \geq \L_o^1 (E_{r,s} \cap V) - \epsilon.
  \end{align*}
  Write $J_k = (y_k, y_k + t_k)$, we have
  \begin{align*}
    \sumk^l f(y_k + t_k) - f(y_k)
     & \geq s \sumk^l t_k                              \\
     & \geq s \L_o^1 (E_{r,s} \cap V) - s \epsilon     \\
     & \geq s \L_o^1 (E_{r,s}) - 2 s \epsilon. \tag{4}
  \end{align*}

  Now compare (3) and (4). Since $f$ is increasing and each $J_k$ lies in
  one of the interval $I_i$, it is not hard to see that
  \[
    \sumi^n f(x_i + h_i) \geq \sumk^l f(y_k + t_k) - f(y_k).
  \]
  It follows that
  \[
    \L_o^1 (E_{r, s}) \leq \frac{2s + r}{s - r} \epsilon.
  \]
  Letting $\epsilon \to 0$ completes the proof for $\L^1 (E_{r, s}) = 0$.
  This implies that $f'_+ (x)$ exists for a.e.\ $x \in [a, b]$. Similarly,
  $f'_-(x)$ exists for a.e.\ $x \in [a, b]$. A previous lemma shows that
  the left and right derivatives are equal a.e. Therefore, $f'(x)$ exists
  a.e.

  It remains to show the left and right derivatives are a.e.\ finite.
  Define
  \begin{align*}
    g(x) = \begin{cases}
             f(x), & x \in [a, b] \\
             f(b), & x \geq b.
           \end{cases}
  \end{align*}
  Define also
  \[
    g_n(x) = \frac{g(x + n^{-1}) - g(x)}{n^{-1}}
  \]
  for $x \in [a, b]$ and $n \in \N$. Then, $g_m (x) \to f'(x)$ for a.e.\
  $x \in [a, b]$. Since $f$ is monotone, it is measurable and thus $g_m$
  is also measurable. By Fatou's lemma,
  \begin{align*}
    \int_a^b f'(x) \d x
     & \leq \liminf_{n \to \infty}  \int_a^b g_n (x) \d x              \\
     & = \liminf_{n \to \infty}
    n \left( \int_a^b g(x + n^{-1}) \d x - \int_a^b g(x) \d x \right)  \\
     & = \liminf_{n \to \infty}
    n \left( \int_a^{b - n^{-1}} g(x + n^{-1}) \d x
    + \int_{b - n^{-1}}^b f(b) - \int_a^b g(x) \d x \right)            \\
     & = \liminf_{n \to \infty}
    n \left( \int_{a + n^{-1}}^b g(t) \d t
    - \int_a^b g(t) \d t  \right) + f(b)                               \\
     & = \liminf_{n \to \infty} f(b) - n \int_a^{a + n^{-1}} f(t) \d t \\
     & \leq f(b) - f(a).
  \end{align*}
  Therefore, $f'$ is integrable, which implies $\{ f' = \infty \}$ has
  measure zero.

  For a general interval $I \subset \R$, $I$ can be written as a countable
  union of bounded interval, the proof for the theorem is complete.
\end{proof}

\begin{remark}
  From the proof above, we see that for monotone function $f : [a, b] \to
    \R$, we have
  \[
    \int_a^b f' \leq f(b) - f(a).
  \]
  There exists a monotone function for which inequality is strict. Indeed,
  let $f : [0, 1] \to [0, 1]$ be the Cantor function. Then, $f' = 0$ a.e.\
  but $f(1) = 1$ and $f(0) = 0$.
\end{remark}

\subsection{Differentiation of BV functions}

\begin{defi}[pointwise variation]
  Let $I \subset \R$ be an interval and $f : I \to \R^N$. The
  \textbf{pointwise variation} of $f$ on the interval $I$ is defined as
  \[
    \Var f = \sup \left\{  \sumi^n \norm{f(x_i) - f(x_{i - 1})} \right\},
  \]
  where the supremum is taken over all partition $x_0 < x_1 < \dots <
  x_n$, where $n \in \N$ and $x_0, \dots, x_n \in I$. Say a function has
  \textbf{bounded pointwise variation} if $\Var f < \infty$.

  The space of all functions $f : I \to \R^N$ with bounded pointwise
  variation is denoted by $BV(I; \R^N)$.
\end{defi}

\begin{prop}
  The following properties holds:
  \begin{enumerate}
    \item Fix $c \in I$, then
          \[
            \norm{f}_{BV} = \Var f + \norm{f(c)}
          \]
          is a norm on $BV(I; \R^N)$. Moreover, $BV(I; \R^N)$ is complete
          under this norm.

    \item Let $c \in I$, then
          \[
            \Var_I f = \Var_{I \cap (-\infty, c]} f
            + \Var_{I \cap [c, \infty)} f.
          \]
  \end{enumerate}
\end{prop}

\begin{proof}
  Exercise.
\end{proof}

\begin{defi}
  Let $f : I \to \R^N$. Say $f$ has \textbf{locally bounded variation} if
  $\Var_{[a, b]} f < \infty$ for all $[a, b] \subset I$. Denote the space
  of functions with bounded local variation as $BV_{\loc} (I ; \R^N)$.
\end{defi}

\begin{eg}
  Let $f : I \to \R^N$ is Lipshitz continuous with Lipshitz constant $L$,
  $[a, b] \subset I$, and $a \leq x_0 < x_1 < \dots < x_n < b$. Then,
  \[
    \sumi^n \norm{f(x_i) - f(x_{i - 1})}
    \leq L \sumi^n (x_i - x_{i - 1})
    \leq L (b - a).
  \]
  Therefore, Lipshitz functions has locally bounded variation.
\end{eg}

\begin{ex}
  Let $f, g \in BV ([a, b])$, then the following holds:
  \begin{enumerate}
    \item $fg \in BV([a, b])$,
    \item Further suppose $g \neq 0$, then $f g^{-1}$ needs not be in $BV
          ([a, b])$. However, if there exists $c > 0$ such that $g(x) \geq
          c$ for all $x \in I$, we have $f g^{-1} \in BV ([a, b])$.
  \end{enumerate}
\end{ex}

\begin{thm}
  Let $I \subset \R$ be an interval and $f : I \to \R$ increasing. Then,
  \[
    \Var f = \sup f - \inf f.
  \]
  In particular, $f \in BV_{\loc}(I)$, and if $f$ is bounded, $f \in
  BV(I)$.
\end{thm}

\begin{proof}
  Let $x_0 < x_1 < \dots < x_n$ be a partition of $I$, then
  \begin{align*}
    \sumi^n \abs{f(x_i) - f(x_{i - 1})}
    = f(x_n) - f(x_0)
    \leq \sup f - \inf f.
  \end{align*}
  Hence, $\Var f \leq \sup f - \inf f$. Conversely, assume $\inf I < \sup
    I$ and take $\inf I < a < b < \sup I$. Then, we have $f(b) - f(a) \leq
    \Var f$. Letting $a \to \inf I$ and $b \to \sup I$ completes the proof.
\end{proof}

\begin{thm}
  Let $I \subset \R$ be an interval and $f : I \to \R^N$. Fix $c \in I$
  and define
  \[
    v(x) = \begin{cases}
      \Var_{[c, x]} f   & (x \geq c), \\
      - \Var_{[x, c]} f & (x < c).
    \end{cases}
  \]
  Then, for every $x, y \in I$ with $x < y$, $\norm{f(x) - f(y)} \leq v(y)
  - v(x)$.

  Moreover, if $f \in BV_{\loc} (I ; \R^N)$, then the left and right limit
  exists for all $x \in I^\circ$, and $f$ is continuous at all but
  countably many points. Also, if $N = 1$, $v \pm f$ are increasing.
\end{thm}

\begin{proof}
  Let $x < y$, then
  \begin{align*}
    \norm{f(y) - f(x)}
    \leq \Var_{[x, y]} f
    =
    \begin{cases}
      \Var_{[c, y]} f - \Var_{[c, x]} f & (c \leq x < y),    \\
      \Var_{[x, c]} f - \Var_{[y, c]} f & (x < y \leq c),    \\
      \Var_{[x, c]} f + \Var_{[c, y]} f & (x \leq c \leq y).
    \end{cases}
  \end{align*}
  It directly follows that $\norm{f(x) - f(y)} \leq v(y) - v(x)$.

  Now assume $f \in BV_{\loc} (I ; \R^N)$, then $v$ is real-valued. Now
  that $v$ is increasing, so $v$ is continuous at all but countably many
  points. It directly follows that $f$ is continuous at all but countably
  many points.

  For $N = 1$, we have $\abs{f(x) - f(y)} \leq v(y) - v(x)$, so $f(x) -
  f(y) \leq v(y) - v(x)$ and $f(y) - f(x) \leq v(y) - v(x)$, showing that
  $v \pm f$ are increasing.
\end{proof}

\begin{cor}
  Let $I \subset \R$ be an interval, then
  \begin{enumerate}
    \item $BV_{\loc} (I)$ is the smallest vector space that contains all
          monotone functions $f : I \to \R$.

    \item $BV (I)$ is the smallest vector space that contains all $f : I
            \to \R$ that is bounded and monotone.
  \end{enumerate}
\end{cor}

\begin{proof}
  \begin{enumerate}
    \item By previous theorems, $BV_{\loc} (I)$ contains all monotone
          functions. Conversely, for any $f \in BV_{\loc} (I)$, we have
          \[
            f = v + (f - v),
          \]
          where both functions are monotone.
    \item Similar.
  \end{enumerate}
\end{proof}

\begin{cor}
  If $f \in BV_{\loc} (I)$, then $f$ is differentiable a.e.
\end{cor}

\subsection{Absolutely continuous functions}

\begin{defi}[absolutely continuous]
  Let $I \subset \R$ be an interval. Say $f : I \to \R^N$ is
  \textbf{absolutely continuous} if for every $\epsilon > 0$, there exists
  $\delta > 0$ such that for any finite collection of intervals
  $\seqi{[a_i, b_i]}^n$ with $[a_i, b_i] \subset I$ and disjoint interior,
  and $\sumi^n (b_i - a_i) \leq \delta$, we have $\sumi^n \norm{f(b_i) -
      f(a_i)} \leq \epsilon$.
\end{defi}

\begin{remark}
  Note that taking $n = 1$ in the previous definition gives uniform
  continuity. Therefore absolute continuity implies uniform continuity.
\end{remark}

\begin{eg}
  Let $f : I \to \R^N$ be Lipshitz continuous with constant $L$, then
  $\sumi^n \norm{f(b_i) - f(a_i)} \leq L \sumi^n (b_i - a_i)$. Therefore,
  Lipshitz continuity implies absolute continuity.
\end{eg}

\begin{thm}
  Let $I \subset \R$ be an interval and $f : I \to \R^N$ absolutely
  continuous. Suppose $f' = 0$ a.e. Then, $f$ is constant.
\end{thm}

\begin{eg}
  The Cantor function has zero derivative a.e.\ but it is not constant, so
  it is not absolutely continuous.
\end{eg}

\begin{proof}
  Let $a < b \in I$. We want to show that $f(a) = f(b)$. Let
  \[
    G = \{ x \in (a, b) : f'(x) = 0 \}.
  \]
  It follows that $\L^1 (G) = b - a$. For any $x \in G$, we have $h^{-1}
  (f(x + h) - f(x)) \to 0$ as $h \to 0$. Given $\epsilon > 0$, there
  exists $h_x > 0$ such that $\norm{h^{-1} (f(x + h) - f(x))} \leq
  \epsilon$ for all $0 < \abs{h} \leq h_x$. Now let
  \[
    \fcal = \{ (x, x + h) : x \in G, 0 < h < h_x \}.
  \]
  By a previous lemma, for any $\delta > 0$, there exists a finite and
  disjoint collection of intervals $\seqi{(x_i, x_i + h_i)}^n \subset
  \fcal$ with
  \begin{align*}
     & \L_o^1 \left( G \cap \cupi^n (x_i, x_i + h_i) \right)
    \geq \L_o^1 (G) - \delta,                                    \\
     & \sumi^n \L^1 ((x_i, x_i + h_i)) \leq \L_o^1 (G) + \delta.
  \end{align*}
  Let $x_1 < \dots < x_n$. It follows that
  \begin{align*}
    \sumi^n h_i
     & \geq \L_o^1 \left( G \cap \cupi^n (x_i, x_i + h_i) \right) \\
     & \geq b - a - \delta.
  \end{align*}
  Therefore, $(a, x_1) \cup (x_1 + h_1, x_2) \cup \dots \cup (x_n + h_n,
    b)$ has measure at most $\delta$. By absolutely continuity, we have
  \[
    \norm{f(x_1) - f(a)} + \norm{f(x_1 + h_1) + f(x_2)}
    + \dots + \norm{f(b) - f(x_n + h_n)} \leq \epsilon.
  \]
  Moreover, $\norm{f(x_i + h_i) - f(x_i)} \leq \epsilon h_i$. It follows
  that
  \begin{align*}
    \norm{f(b) - f(a)}
     & \leq \epsilon + \sumi^n \epsilon h_i       \\
     & \leq \epsilon + \epsilon (b - a) + \delta.
  \end{align*}
  Letting $\delta \to 0$ and $\epsilon \to 0$ shows $f(b) = f(a)$.
\end{proof}

\begin{thm}
  Let $I \subset \R$ be an interval and $f : I \to \R^N$ be absoultely
  continuous. Then, $f \in BV_\loc (I ; \R^N)$. In particular, $f$ is
  differentiable a.e.
\end{thm}

\begin{proof}

  Let $[a, b] \in I$. Let $\delta$ be in the definition of absolute
  continuity when applying it to $f$ with $\epsilon = 1$. Divide interval
  $[a, b]$ into segments of length $\delta$. Call the segments
  $\seqi{[x_{i - 1}, x_i]}^n$, where $n \leq \delta^{-1}(b - a) + 1 <
  \infty$. It follows that
  \begin{align*}
    \Var_{[a, b]} f
    = \sumi^n \Var_{[x_{i - 1}, x_i]} f
    \leq n
    < \infty,
  \end{align*}
  completing the proof.
\end{proof}

\subsection{The fundamental theorem of calculus}

\begin{thm}[FTC for Lipshitz functions]
  Let $I \subset \R$ be an interval and $f : I \to \R$ be Lipshitz. Then,
  \[
    f(b) - f(a) = \int_a^b f'(x) \d x.
  \]
\end{thm}

\begin{proof}
  Note that $f$ being Lipshitz implies $f \in BV_\loc (I)$, so $f$ is
  differentiable a.e. Define
  \[
    g(x) = \begin{cases}
      f(x) & (a \leq x \leq b), \\
      f(b) & (x \geq b).
    \end{cases}
  \]
  Note that $g$ is still Lipshitz and let $L$ be the Lipshitz constant. We
  have
  \begin{align*}
    \lim_{n \to \infty} \frac{g(x + n^{-1}) - g(x)}{n^{-1}}
    = g'(x) \text{ a.e.}
  \end{align*}
  Since $g$ is continuous, $g'$ is measurable. Moreover, $g$ being
  Lipshitz implies $\abs{g'} \leq L$. It then follows from Dominated
  Convergence Theorem that
  \[
    \lim_{n \to \infty} \int_a^b \frac{g(x + n^{-1}) - g(x)}{n^{-1}} \d x
    = \int_a^b g'(x) \d x.
  \]
  Through similar computation as in the Lebesgue differentiation theorem,
  for LHS we have
  \begin{align*}
    \lim_{n \to \infty} \int_a^b \frac{g(x + n^{-1}) - g(x)}{n^{-1}} \d x
     & = \lim_{n \to \infty} f(b) - n \int_a^{a + n^{-1}} f(t) \d t \\
     & = f(b) - f(a),
  \end{align*}
  where the last equality is by continuity of $f$.
\end{proof}

\begin{remark}
  Comparing this with the proof of Lebesgue differentiation theorem for
  monotone function, we know $f$ is continuous and Lipshitz. Lipshitz
  continuity gives us bounded derivative and the ability to use Dominated
  Convergence Theorem, and continuity let us know the final term is
  $f(a)$.
\end{remark}

\begin{thm}[FTC 1]
  Let $g : [a, b] \to \R$ integrable. Let $f(x) = \int_a^x g(t) \d t$.
  Then $f$ is absolutely continuous, and $f'(x) = g(x)$ for a.e.\ $x \in
    [a, b]$.
\end{thm}

To show this we need a few lemmas.

\begin{lemma}
  Let $E \subset \R^N$ be Lebesgue measurable, $g : E \to \R$ be
  integrable. Then for all $\epsilon > 0$, there exists $\delta > 0$ such
  that $F \subset E$ with $\L^N (F) \leq \delta$ implies $\int_F \abs{g}
    \d x \leq \epsilon$.
\end{lemma}

\begin{lemma}
  Let $g : [a, b] \to \R$ integrable such that $\int_a^x g(t) \d t = 0$
  for all $x \in [a, b]$. Then, $g = 0$ a.e.
\end{lemma}

\begin{proof}
  Exercise.
\end{proof}

Now we begin the proof for the first fundamental theorem of calculus.

\begin{proof}
  First we show $f \in AC$. Fix $\epsilon > 0$, and let $\delta$ be as in
  the first lemma. Suppose $\seqi{(a_i, b_i)}^n$ are disjoint sub-interval
  of $[a, b]$ with $\sumi^n (b_i - a_i) \leq \delta$. Let $U = \cupi^n
    (a_i, b_i)$. It follows that
  \begin{align*}
    \sumi^n \abs{f(b_i) - f(a_i)}
    = \sumi^n \abs{\int_{a_i}^{b_i} g(t) \d t}
    \leq \sumi^n \int_{a_i}^{b_i} \abs{g(t)} \d t
    = \int_U \abs{g(t)} \d t
    \leq \epsilon,
  \end{align*}
  as desired. Since $f \in AC$, it is differentiation a.e.

  It remains to show that $f'(x) = g(x)$ a.e. First assume $g$ is bounded
  with $\abs{g} \leq M$. In this case, $f$ is Lipshitz. Then for any $c
  \in [a, b]$, we have $f(c) - f(a) = \int_a^c f'(t) \d t$ by FTC for
  Lipshitz functions. It follows that
  \[
    \int_a^c g(t) \d t = \int_a^c f'(t) \d t.
  \]
  By the second lemma above, we have $f' = g$ a.e.

  For the next step assume $g$ is non-negative. Define
  \[
    g_n(t) = \begin{cases}
      g(t) & (g(t) \leq n), \\
      0    & (g(t) > n).
    \end{cases}
  \]
  We then have $0 \leq g_n \leq n$. Moreover,
  \begin{align*}
    f(x)
    = \int_a^x g(t) \d t
    = \int_a^x g_n(t) \d t + \int_a^x (g(t) - g_n(t)) \d t.
  \end{align*}

  Write $G_n(x) = \int_a^x g_n(t) \d t$ and $H_n = \int_a^x (g(t) -
  g_n(t)) \d t$. By the previous step, we know that $G_n' (x) = g_n(x)$
  a.e. We also know that $g \geq g_n$, so $H_n$ is increasing. Then, $H_n$
  is differentiable a.e.\ and $f' = g_n + H_n' \geq g_n$ a.e. Letting $n
  \to \infty$ gives $f' \geq g$ a.e.

  To get equality, we take integral and obtain $\int_a^b f'(x) \d x \geq
  \int_a^b g(x) \d x$. Since $g \geq 0$, $f$ is also increasing. It
  follows that $\int_a^b f'(x) \d x \leq f(b) - f(a)$. However $\int_a^b
  g(x) \d x = f(b) - f(a)$. This shows that $\int_a^b (f' - g) = 0$. With
  $f' \geq g$, this implies $f' = g$ a.e.

  To complete the proof for general $g$, consider $g = g^+ - g^-$.
\end{proof}

\begin{thm}[FTC 2]
  Let $f : [a, b] \to \R$, then $f \in AC$ iff the following holds:
  \begin{enumerate}
    \item $f$ is differentiable a.e.
    \item $f'$ is Lebesgue integrable.
    \item $f'(x) - f'(x_0) = \int_{x_0}^x f'(t) \d t$ for all
          $x, x_0 \in [a, b]$.
  \end{enumerate}
\end{thm}

\begin{proof}
  Suppose the three given conditions holds. Let $F(x) = \int_a^x f'(t) \d
    t$. By the previous theorem (FTC 1), we know $F \in AC$. Meanwhile,
  condition 3 implies $f(x) = f(a) + F(x)$. This shows that $f \in AC$.

  Conversely, suppose now that $f \in AC$. We have shown before that $f$
  is differentiable a.e. Moreover, $f$ is the difference of two increasing
  functions, while the derivative of increasing functions is integrable.
  This shows that $f'$ is integrable.

  Again let $F(x) = \int_a^x f'(t) \d t$. By FTC 1, $F \in AC$ and $F' =
  f'$ a.e. Moreover $f - F \in AC$ and $(f - F)' = 0$ a.e. By previous
  theorem, this implies $f - F = c$ a.e. for some constant $c$. Therefore,
  $f(x) = \int_a^x f'(x) \d t + c$ and the proof is complete.
\end{proof}

\section{Change of variables}

\subsection{Mollifiers}

\begin{defi}[support]
  Let $(X, d)$ a metric space, $E \subset X$, and $f : E \to \R$. The
  \textbf{support} of $f$ is
  \[
    \supp f = \bar{\{ x \in E : f(x) \neq 0 \}}.
  \]
  Write the set of continuous function with compact support as $C_c (X)$.
  Also for open set $U \subset X$, write $C_c^k (U) = C_c(U) \cap C^k
  (U)$.
\end{defi}

\begin{ex}
  Define
  \[
    \phi(x) = \begin{cases}
      e^{(\norm{x}^2 - 1)^{-1}} & (\norm{x} < 1),    \\
      0                         & (\norm{x} \geq 1).
    \end{cases}
  \]
  Show that $\phi \in C_c^\infty(\R^N)$.
\end{ex}

\begin{defi}[mollifier]
  Let $\phi \in C_c(\R^N)$ with $\supp \phi \subset B[0, 1]$ and $\int \phi
    = 1$. For any $\epsilon > 0$, define
  \[
    \phi_\epsilon (x) = \epsilon^{-N} \phi(\epsilon^{-1} x).
  \]
  These functions are called \textbf{mollifiers}. Note that $\supp
  \phi_\epsilon \subset B[0, \epsilon]$.
\end{defi}

\begin{ex}
  Recall that for $\lambda > 0$, we have $\L^N (\lambda E)
    = \lambda^N \L^N (E)$. Show that $\int \phi_\epsilon = 1$.
\end{ex}

\begin{defi}[mollification]
  Let $E \subset \R^N$ be measurable, $f : E \to \R$ integrable. For $x \in
    \R^N$, define
  \[
    f_\epsilon (x) = \int_E \phi_\epsilon \left( x - y \right)
    f (y) \d y.
  \]
  The function $f_\epsilon$ is the \textbf{mollification} of $f$.
\end{defi}

\begin{thm}
  Let $E \subset \R^N$ be measurable, $f : E \to \R$ integrable, then
  $f_\epsilon \in C^\infty (\R^N)$, and
  \[
    \frac{\dd^\alpha f_\epsilon}{\dd x^\alpha} (x)
    = \int_E \frac{\dd^\alpha \phi_\epsilon}{\dd x^\alpha} (x - y) f(y) \d y.
  \]
\end{thm}

\begin{proof}
  First show $f_\epsilon$ is continuous. Fix $x \in \R^N$ and let
  $\seqinfn{x_n}$ be a sequence with $x_n \to x$. Note that
  \[
    \abs{\phi_\epsilon (x_n - y) f(y)}
    \leq \sup \abs{\phi_\epsilon} \abs{f(y)},
  \]
  which is integrable. It then follows from Dominated Convergence Theorem
  that $f_\epsilon (x_n) = f(x)$ as $n \to \infty$.

  For the paritial derivative, fix $i \in [N]$. We have
  \begin{align*}
     & \frac{f_\epsilon(x + t e_i) - f_\epsilon (x)}{t}
    - \int_E \frac{\dd \phi_\epsilon}{\dd x_i} (x - y) f(y) \d y                  \\
     & = \int_E f(y) \left[
                       \frac{\phi_\epsilon (x + t e_i) - \phi_\epsilon (x - y)}{t}
                       - \frac{\dd \phi_\epsilon}{\dd x_i} (x - y) \right]
    \d y.
  \end{align*}
  By mean value theorem, we have
  \begin{align*}
    \phi_\epsilon (x + t e_i - y)
    - \phi_\epsilon (x - y)
    = t \frac{\dd \phi_\epsilon}{\dd x_i} (x + \theta e_i - y),
  \end{align*}
  for some $\abs{\theta} \leq \abs{t}$. Hence the absolute value of the
  integrand is at most $2 \abs{f(y)} \sup \abs{\dd_i \phi_\epsilon}$,
  which is integrable. Then Dominated Convergence Theorem applies and
  shows the formula for the paritial derivative.

  By induction, we can show that $f_\epsilon \in C^\infty (\R^N)$ and
  $\dd_\alpha f_\epsilon$ is as desired.
\end{proof}

Now we discuss some properties of the mollification.

\begin{thm}
  Let $E \subset \R^N$ be measurable and $f : E \to \R^N$ be integrable.
  Assume $f$ is continuous at $x \in E^\circ$. Then, $f_\epsilon (x) \to
    f(x)$ as $\epsilon \to 0^+$.
\end{thm}

\begin{proof}

  Since $x \in E^\circ$, there exists $B(x_0, r) \subset E$. For $\epsilon
  < r$, we have
  \begin{align*}
    f_\epsilon (x) - f(x)
     & = \int_{E \cap B(x, \epsilon)} \phi_\epsilon (x - y) f(y) \d y - f(x) \\
     & = \int_{\R^N} \phi_\epsilon (x - y) f(y) \d y - f (x)                 \\
     & = \int_{\R^N} \phi_\epsilon(x - y) (f(y) - f(x)) \d y.
  \end{align*}
  Fix $\eta > 0$. There exists $\delta > 0$ such that $\norm{y - x} <
    \delta$ implies $\abs{f(y) - f(x)} < \eta$. Then when $\epsilon <
    \delta$, we have
  \begin{align*}
    \abs{f_\epsilon (x) - f(x)}
    \leq \eta \int_{\R^N} \phi_\epsilon (x - y) \d y
    = \eta.
  \end{align*}
  The proof is then complete.
\end{proof}

\end{document}